\section{Tests Beyond the Benchmark}
\label{sec:beyond}
\subsection{Budget Scaling and Feasibility-Preserving Initialization}
\label{sec:feasinit}
\label{sec:budget}
%% generated by mpce_results.py -- do not edit by hand
\begin{table}[!t]
\centering
\caption{Six Largest Benchmark Cases: Mean Wake Loss (\%) and Feasible Runs (\%) at 6,030 Calls with Random and Feasibility-Preserving Initialization, and Average Rank at 6,030, 30,030 and 120,030 Calls (Random Initialization; 30 Seeds)}
\label{tab:feasbudget}
\scriptsize\setlength{\tabcolsep}{2.4pt}
\begin{tabular}{lccccccc}
\toprule
& \multicolumn{2}{c}{Random init.} & \multicolumn{2}{c}{Feasible init.} & \multicolumn{3}{c}{Avg.\ rank} \\
\cmidrule(lr){2-3}\cmidrule(lr){4-5}\cmidrule(lr){6-8}
Method & Loss & Feas. & Loss & Feas. & 6,030 & 30,030 & 120,030 \\
\midrule
PSO-VNS & 4.084 & 100.0 & 3.938 & 100.0 & 1.50 & 1.17 & 2.17 \\
SSA-VNS & 4.678 & 95.0 & \TBD{} & \TBD{} & 4.33 & 4.17 & 4.17 \\
LX-SSA-VNS & 4.802 & 98.9 & \TBD{} & \TBD{} & 5.17 & 5.67 & 5.83 \\
RS-VNS & 4.830 & 93.3 & \TBD{} & \TBD{} & 4.67 & 6.33 & 4.67 \\
VNS & 4.950 & 100.0 & \TBD{} & \TBD{} & 6.50 & 4.83 & 3.17 \\
PSO & 4.181 & 100.0 & \TBD{} & \TBD{} & 1.83 & 2.50 & 5.83 \\
SSA & 4.828 & 80.0 & \TBD{} & \TBD{} & 8.33 & 8.67 & 8.83 \\
LX-SSA & 5.379 & 77.8 & \TBD{} & \TBD{} & 8.67 & 9.83 & 10.00 \\
DE & -- & 1.1 & \TBD{} & \TBD{} & 10.00 & 5.17 & 3.33 \\
MS-SLSQP & 4.522 & 99.4 & \TBD{} & \TBD{} & 4.00 & 6.67 & 7.00 \\
\bottomrule
\multicolumn{8}{l}{\TBD{pending: feas}}
\end{tabular}
\end{table}


\emph{Budget.} On the six largest benchmark cases, \NBudgetPhrase{} (6,030, 30,030 and 120,030 evaluations: \NBudRankPSOVNSSixK, \NBudRankPSOVNSThirtyK{} and \NBudRankPSOVNSOneTwentyK; Table~\ref{tab:feasbudget}; Fig.~\ref{S-fig:S-budget}); it ranks first in \NBudFirstPSOVNSSixK, \NBudFirstPSOVNSThirtyK{} and \NBudFirstPSOVNSOneTwentyK{} of the six cases (exploratory; rank differences not tested). PSO's mean wake loss minus that of PSO-VNS is \NBudGainPSOSixK, \NBudGainPSOThirtyK{} and \NBudGainPSOOneTwentyK{} percentage points (per case: Table~\ref{S-tab:budget-cases}). The order of the other methods depends on the budget. PSO falls back from an average rank of \NBudRankPSOSixK{} to \NBudRankPSOOneTwentyK, so that at 120,030 evaluations SSA-VNS ranks ahead of it (\NBudRankSSAVNSOneTwentyK), while VNS (\NBudRankVNSSixK{} to \NBudRankVNSOneTwentyK) and DE move up. At 120,030 evaluations, VNS, SSA-VNS, RS-VNS, LX-SSA-VNS and DE all come within \NBudCloseGapOneTwentyK~percentage points of the mean wake loss of PSO-VNS (\NBudLossPSOVNSOneTwentyK\%).
% CHECK-FINAL [C26]: feasbudget.rank.6030: PSOBV best (column 6,030; the 30,030 / 120,030 statements are generated by \NBudgetPhrase)
% CHECK-FINAL [C29]: PSO-VNS best average rank at all three budgets
% CHECK-FINAL [C37] (proposed in R5-16, to be defined in mpce_check_final.py): feasbudget.rank PSOC[120030] > PSOC[6030] + 2; BVNS and DE [120030] < [6030] - 2; SSABV[120030] < PSOC[120030]; feasbudget.mean_loss[120030] of BVNS, SSABV, RSVNS, LXBV, DE minus PSOBV < 0.2 pp
% verified 2026-09-28 from mpce_summary.json (feasbudget.rank, feasbudget.mean_loss): PSOC rank 1.83 -> 5.83, BVNS 6.50 -> 3.17, DE 10.00 -> 3.33, SSABV 4.17 < PSOC 5.83 at 120,030; losses at 120,030: PSOBV 3.145, BVNS 3.174, SSABV 3.228, RSVNS 3.255, LXBV 3.306, DE 3.307 (max gap 0.163 pp)

\emph{Feasible starts.} With feasibility-preserving initialization, \NFeasInitPhrase. \NFbFeasBestRank{} alone is then best in mean wake loss and average rank; PSO-VNS follows (rank \NFbFeasRankPSOVNS, loss \NFbFeasLossPSOVNS\%; \NFbRandLossPSOVNS\% with random starts; Table~\ref{S-tab:feasinit-cases}). PSO and DE never improve on the best initial layout (\NDFsStoredRunsPSO{} runs each).
\makeatletter\@ifundefined{proposition}{\newtheorem{proposition}{Proposition}}{}\makeatother
\begin{proposition}\label{prop:feasible-start}
From feasible starts with $\mathbf V^i=\mathbf 0$ and $\mathbf P^i=\mathbf X^i$, the particle holding $\mathbf G$ keeps $\mathbf V=\mathbf 0$ until $\mathbf G$ changes, and a PSO or DE trial replaces $\mathbf G$ only if it is feasible and has a lower wake loss.
\end{proposition}
In our runs, no trial passes this test (proofs: Section~\ref{S-sec:S-th-feas}; data: Table~\ref{S-tab:D-feasstart}): apart from the re-evaluations of $\mathbf G$, all PSO candidates and all \NDFsCandDE{} DE trials are infeasible, and no personal best ever changes. PSO moves each coordinate by an independent random weight, with overshoot, toward the same coordinate of $\mathbf G$, and DE takes most coordinates from a combination of three other layouts; for distinct dense layouts, both make turbines collide (median minimum spacing at the first move \NDFsFirstMinSpPSO~m, limit \NDSmin~m), and most trials also place a turbine outside the site. Relabelling the turbines to match those of $\mathbf G$ does not help (\NDFsMixMatched\% of first PSO moves feasible). The stall is observed, not implied; by the proposition, only a feasible, better trial can end it. VNS, which moves one coordinate at a time, is not affected. From random starts, PSO may thus help PSO-VNS mainly by reaching feasibility rather than by a better wake-loss search, and the advantage of PSO-VNS depends on the initialization.
% CHECK-DIAG [D13]: PSO and DE never improve on the best initial layout: instrumented runs (0 of \NDFsRunsPSO{} each) and all stored runs (\NDFsStoredMovedPSO{} / \NDFsStoredMovedDE{} of \NDFsStoredRunsPSO{} final != best initial)
% CHECK-DIAG [D14]: every feasible PSO candidate is the G-particle (V = 0) re-evaluating G (\NDFsOneFeasIter{} of \NDFsIter{} iterations with exactly one feasible candidate; \NDFsFrozen{} of \NDFsRunsPSO{} runs frozen); no pbest update; DE 0 feasible of \NDFsCandDE{} trials, 0 accepted; first-move median min spacing \NDFsFirstMinSpPSO{} m (PSO), \NDFsFirstMinSpDE{} m (DE); infeasible with a turbine outside: \NDFsInfBothPSO{} % (PSO), \NDFsInfBothDE{} % (DE) ("most")
% CHECK-DIAG [D15]: relabelling does not help: first PSO move feasible in \NDFsMixRaw{} % (original labels) and \NDFsMixMatched{} % (relabelled) of draws (< 1 %)
% CHECK-DIAG [D16]: the feasible initial layouts are distinct ("distinct dense layouts"): matched RMS distance to G median \NDFsInitRms{} m, min \NDFsInitRmsMin{} m (> 50 m)
% CHECK-THEORY: Proposition prop:feasible-start = Lemma S-lem:S-feas-improve + Proposition S-prop:S-frozen(a) (proofs in optA/supp_theory.tex, Section S-th-feas); pure code facts, no T-check needed; "DE takes most coordinates from the mutant" = Proposition S-prop:S-de, mean 1+CR(n-1) = 27.1 of 30 for N = 15 [T11]. The stall itself is an observation (D13/D14), NOT a consequence of the proposition.
% The former label-symmetry explanation ("as the turbine labels are arbitrary ... the global best of PSO does not move") was removed in Phase 6: relabelling does not restore feasibility (D15).
% verified 2026-09-28 from mpce_feas_s0of1.csv: the convergence curves of PSO and DE show no improvement after the initial population in any of the 210 feasible-initialization runs (identical final layouts for both).
% verified 2026-09-28 from mpce_summary.json (feasbudget.feasible, rank_feasible_init): BVNS lowest loss 3.810 (PSOBV 3.938) and rank 1.00 (PSOBV 2.83); PSOC and DE loss identical (5.5488)
% CHECK-FINAL [C28]: feasible starts: VNS alone first, PSO-VNS second
% CHECK-FINAL [C38] (proposed in R5-17, to be defined in mpce_check_final.py): argmin feasbudget.feasible.*.loss == argmin rank_feasible_init == BVNS; |loss PSOC - loss DE| < 1e-9 (neither moves from the best initial layout)

\subsection{Horns Rev~1 16-Turbine Block}
\label{sec:hornsrev-site}
%% generated by mpce_results.py -- do not edit by hand
\begin{table}[!t]
\centering
\caption{Horns Rev~1 16-Turbine Block (Wake-Free AEP 149.26 GWh/yr): Mean AEP (GWh/yr) over the Feasible Runs of 30 Seeds at 6,030 Calls, Feasible Runs, Holm-Adjusted Wilcoxon $p$ of PSO-VNS vs.\ Each Method, and Mean AEP Loss (\%) at 6,030 Calls with Random (R) and Feasibility-Preserving (F) Initialization and at 30,030 and 120,030 Calls (R; 10 Seeds at 120,030); Superscript: Feasible Runs When Not All Runs Are Feasible}
\label{tab:hr-site}
\scriptsize\setlength{\tabcolsep}{2.2pt}
\begin{tabular}{lccccccc}
\toprule
& \multicolumn{3}{c}{6,030 calls, R} & \multicolumn{4}{c}{Loss (\%)} \\
\cmidrule(lr){2-4}\cmidrule(lr){5-8}
Layout / method & AEP & Feas. & $p_{\rm Holm}$ & 6k R & 6k F & 30k & 120k \\
\midrule
Installed layout & 139.51 & -- & -- & 6.53 & -- & -- & -- \\
\midrule
PSO-VNS & 138.73 & 30/30 & -- & 7.05 & 6.96 & 6.46 & \TBD{} \\
SSA-VNS & 137.98 & 30/30 & $6.7\times10^{-4}$ & 7.56 & \TBD{} & \TBD{} & \TBD{} \\
LX-SSA-VNS & 137.77 & 30/30 & $6.3\times10^{-4}$ & 7.70 & \TBD{} & \TBD{} & \TBD{} \\
RS-VNS & 137.96 & 28/30 & $2.2\times10^{-4}$ & 7.57$^{28}$ & \TBD{} & \TBD{} & \TBD{} \\
VNS & 137.67 & 30/30 & $2.2\times10^{-4}$ & 7.77 & \TBD{} & \TBD{} & \TBD{} \\
PSO & 138.06 & 19/30 & $1.4\times10^{-5}$ & 7.51$^{19}$ & \TBD{} & \TBD{} & \TBD{} \\
SSA & 136.66 & 30/30 & $2.8\times10^{-7}$ & 8.44 & \TBD{} & \TBD{} & \TBD{} \\
LX-SSA & 136.46 & 28/30 & $1.5\times10^{-7}$ & 8.58$^{28}$ & \TBD{} & \TBD{} & \TBD{} \\
DE & -- & 0/30 & $1.7\times10^{-8}$ & -- & \TBD{} & \TBD{} & \TBD{} \\
MS-SLSQP & 136.80 & 30/30 & $1.5\times10^{-7}$ & 8.35 & \TBD{} & \TBD{} & \TBD{} \\
\bottomrule
\end{tabular}
\end{table}


The model (Section~\ref{S-sec:S-hrmodel}) uses the Jensen wake with a hub-center in-wake test and 5\textdegree{} direction and 1-m/s speed bins; the boundary violation in~\eqref{eq:penalty} is the distance (m) outside the parallelogram. For the 80-turbine farm, our AEP agrees with PyWake~\cite{PyWake} within \NHRPyWakeDiff\%. The installed layout, a regular $7D$ grid, reaches \NHRInstalled~GWh/yr (wake loss \NHRInstalledLoss\%; Table~\ref{tab:hr-site}; layouts: Fig.~\ref{S-fig:S-hr16}). At 6,030 evaluations, PSO-VNS is feasible in \NHRFeasPSOVNS{} runs, PSO in only \NHRFeasPSO{} (\NHRFeasFeasInitPSO{} with feasible starts); PSO-VNS has the highest mean AEP, \NHRMeanPSOVNS~GWh/yr (SD \NHRSDPSOVNS); in the seed-paired run-level test (infeasible runs ranked last) it is significantly better than every other method ($p_{\rm Holm}\le\NHRMaxP$). Its best run reaches \NHRBestPSOVNS~GWh/yr, and \NHRAbovePSOVNS{} of its 30 runs exceed the installed layout. At 30,030 evaluations, PSO-VNS still has the highest mean AEP, \NHRMeanThirtyKPSOVNS~GWh/yr (loss \NHRLossThirtyKPSOVNS\%), and \NHRAboveThirtyKPSOVNS{} of its 30 runs exceed the installed layout. At 120,030 evaluations (10 seeds), its loss of \NHRLossOneTwentyKPSOVNS\% is no longer the lowest; \NHRBestHundredTwentyK{} has the highest mean AEP. The optimized layouts may use the benchmark spacing $4D$, whereas the installed grid keeps $7D$ (no run with $7D$ was made); the block is evaluated without the wakes of the other 64 turbines; and cables, foundations, loads and navigation, which also shaped the installed grid, are not modeled. The comparison therefore tests the optimizers; it does not show that the installed layout could be improved.
% Horns Rev numbers: all \NHR... macros come from the rerun with the corrected direction binning (hrfix, 970 runs); sentences rewritten by the lead 2026-09-28. The mean at 6,030 and at 30,030 is deliberately NOT compared with the installed layout (R4-1, R4-11, R5-2).
% CHECK-FINAL [C19]: hr16 (6,030, random): PSOBV highest mean AEP, all p_holm < 0.05 (best run vs installed stated via macros only)
% CHECK-FINAL [C20]: hr16.methods.PSOC.feasible < runs ("only"); Phase 6: moved next to the PSO-VNS count to state the feasibility-reliability advantage (30/30 vs 19/30; W1 note C52: the robustness advantage is feasibility, not spread)
% CHECK-FINAL [C25]: hr16.loss_by_setting.PSOBV: 30030R and 120030R loss < 6030R ("falls")
% CHECK-FINAL [C39] (proposed in R5-18, to be defined in mpce_check_final.py): hr16.loss_by_setting: some method has 120030R loss < PSOBV 120030R loss ("no longer the lowest"), and PSOBV has the lowest 6030R and 30030R loss
% (C30 no longer asserted here: the 30,030 mean is reported without comparison with the installed layout.)
% verified 2026-09-28 from mpce_tab_hr16.tex (pre-rerun data): at 120,030 BVNS 6.12, RSVNS 6.24, DE 6.34 < PSOBV 6.45; PSOBV lowest at 6k R (7.05) and 30k (6.46)

\subsection{IEA37 Case Study~1}
\label{sec:iea37}
%% generated by mpce_results.py -- do not edit by hand
\begin{table}[!t]
\centering
\caption{IEA37 Case Study~1, 16- and 36-Turbine Scenarios: AEP in MWh of the Best Run of Each Method at 6,030 and 30,030 Calls (30 Seeds), Compared with the Example Layout and the Best Feasible Published Layout~\cite{Baker2019,IEA37repo}}
\label{tab:iea37}
\scriptsize\setlength{\tabcolsep}{2.5pt}
\begin{tabular}{lcccc}
\toprule
& \multicolumn{2}{c}{16 turbines ($r=1300$~m)} & \multicolumn{2}{c}{36 turbines ($r=2000$~m)} \\
\cmidrule(lr){2-3}\cmidrule(lr){4-5}
Method / source & 6,030 & 30,030 & 6,030 & 30,030 \\
\midrule
Example layout & \multicolumn{2}{c}{366{,}941.6} & \multicolumn{2}{c}{737{,}883.1} \\
Best feasible published$^{a}$ & \multicolumn{2}{c}{418{,}924.4} & \multicolumn{2}{c}{863{,}676.3} \\
\midrule
PSO-VNS & \TBD{} & \TBD{} & \TBD{} & \TBD{} \\
PSO & 403{,}753.1 & 412{,}702.4 & 787{,}988.7 & 830{,}256.4 \\
SSA-VNS & 403{,}251.3 & 406{,}924.6 & 766{,}803.3 & 819{,}930.5 \\
SSA & 395{,}150.2 & 401{,}459.3 & 750{,}707.4 & 802{,}538.2 \\
LX-SSA & 393{,}284.6 & 398{,}770.0 & 746{,}301.3 & 789{,}311.3 \\
DE & 357{,}391.3 & 397{,}402.7 & -- & -- \\
VNS & 400{,}169.3 & 404{,}076.0 & 785{,}791.3 & 824{,}230.2 \\
MS-SLSQP & 405{,}433.8 & 410{,}349.5 & 748{,}571.1 & 833{,}365.0 \\
\bottomrule
\multicolumn{5}{p{0.92\columnwidth}}{$^{a}$Participant~4 in both scenarios (SNOPT with wake expansion continuation \TBD{verify against Baker et al.\ (2019)}). A higher submitted AEP places turbines up to 3.5~m outside the boundary and is not counted.}
\end{tabular}
\end{table}

% Fig. fig:layouts (figures_mpce/layouts_iea37_hr16.pdf) moved to the supplement for the page budget (D10); proposed caption there:
% "Best layouts. Left and middle: IEA37 Case Study~1, 16 and 36 turbines, best feasible published layout and best PSO-VNS layout at 30,030 evaluations. Right: Horns Rev~1 16-turbine block, installed layout ($7D$ grid) and best PSO-VNS layout at 6,030 evaluations (minimum spacing $4D$)."

In IEA37 Case Study~1 the participants chose their own optimizers and budgets~\cite{Baker2019}. None of our methods reaches the best feasible published layouts (Table~\ref{tab:iea37}; layouts: Section~\ref{S-sec:S-iea37}). At 30,030 evaluations, the relative AEP difference to them is \NIEAGapSixteenThirtyK\% (16 turbines) and \NIEAGapThirtySixThirtyK\% (36 turbines) for the best PSO-VNS layout and \NIEAMeanGapSixteenThirtyK\% and \NIEAMeanGapThirtySixThirtyK\% for the mean over the 30 runs; the best PSO-VNS layout would rank \NIEARankSixteenThirtyK{} among \NIEAPubFeasSixteen{} and \NIEARankThirtySixThirtyK{} among \NIEAPubFeasThirtySix{} feasible submissions. The gap is larger for 36 turbines, where the best of our layouts comes from \NIEABestOfThirtySixThirtyK. As the participants' budgets are not reported comparably, the size of the gap is indicative only. The 64-turbine scenario was not run, and, like the published layouts, ours are optimized for 16 discrete directions and one wind speed.
% verified 2026-09-28 from mpce_numbers.tex / mpce_summary.json (iea37): best of all methods 16T/30k 413,331.9 (PSO-VNS) < 418,924.4 and 36T/30k 833,365.0 (MS-SLSQP) < 863,676.3; gap 16/30k -1.33 % vs 36/30k -4.20 % ("larger for 36 turbines"); rank macros refer to the best PSO-VNS layout
% CHECK-FINAL [C40] (proposed in R5-19, to be defined in mpce_check_final.py): iea37 16T and 36T: best of all methods at 30,030 < best feasible published; |gap36| > |gap16|

\subsection{Computational Cost and Published Results}
\label{sec:literature}
\label{sec:runtime}
Per evaluation, the metaheuristics need \NMsEvalMin--\NMsEvalMax~ms and MS-SLSQP \NMsEvalSLSQP~ms, about \NSLSQPSlowdown{} times as long as PSO (indicative; Table~\ref{S-tab:cost}); one 6,030-evaluation run of PSO-VNS takes \NSecPerRunMin--\NSecPerRunMax~s, depending on $N$. Published Kusiak--Song results~\cite{Kusiak2010,Eroglu2012,Eroglu2013,Bansal2017} are discussed in Section~\ref{S-sec:S-literature}; they are not compared numerically because budgets, constraint handling and reporting differ.
